\documentclass{amsart}

% 数学相关宏包
\usepackage{amsmath,mathtools,amsthm,amssymb}
\usepackage{bbm}
\usepackage{mathrsfs}
\usepackage{centernot}

% 图形和绘图
\usepackage{graphicx,float}
\usepackage{tikz,tikz-cd}
\usepackage{pgfplots}
\usepackage[framemethod=TikZ]{mdframed}
\usetikzlibrary{decorations.pathmorphing,positioning,arrows,automata,calc,shapes,patterns}
\pgfplotsset{compat=1.18}
\usepgfplotslibrary{statistics}

% 表格和列表
\usepackage{booktabs,multirow,colortbl}
\usepackage{enumitem}
\setlist[itemize,enumerate,description]{noitemsep}
\setlistdepth{9}

% 算法
\usepackage{algorithm,algorithmic}

% 字体和颜色
\usepackage{xcolor,listings}
\usepackage{xeCJK}
\setCJKmainfont{Noto Serif CJK SC}

% 页面和链接
\usepackage{fancyhdr,caption}
\usepackage{hyperref}
\usepackage{cleveref}

% 工具宏包
\usepackage{etoolbox,letltxmacro,docmute}
\usepackage{gensymb}

% 自定义设置
\renewcommand{\arraystretch}{1.5}
\let\originalmathbb\mathbb
\renewcommand{\mathbb}[1]{%
    \ifstrequal{#1}{1}{\mathbbm{#1}}{\originalmathbb{#1}}%
}

% 颜色定义
\definecolor{lightblue}{RGB}{173,216,230}
\definecolor{lightgreen}{RGB}{144,238,144}
\definecolor{lightyellow}{RGB}{255,255,224}
\definecolor{lightgray}{RGB}{211,211,211}

% 超链接设置
\hypersetup{
    colorlinks=true,
    linkcolor=blue,
    filecolor=magenta,
    urlcolor=cyan,
    pdftitle={\@title},
    pdfauthor={\@author},
    pdfsubject={数学笔记},
    pdfkeywords={数学},
    bookmarksnumbered=true,
    bookmarksopen=true,
    bookmarksopenlevel=6,
    pdfstartview=FitH,
    pdfpagemode=UseOutlines,
    pdfborder={0 0 0}
}

% 定理环境
\theoremstyle{plain}
\newtheorem{theorem}{Theorem}[section]
\newtheorem{lemma}[theorem]{Lemma}
\newtheorem{corollary}[theorem]{Corollary}
\newtheorem{proposition}[theorem]{Proposition}

\theoremstyle{definition}
\newtheorem{definition}[theorem]{Definition}
\newtheorem{example}[theorem]{Example}
\newtheorem{exercise}[theorem]{Exercise}
\newtheorem{problem}[theorem]{Problem}

\theoremstyle{remark}
\newtheorem{remark}[theorem]{Remark}
\newtheorem{note}[theorem]{Note}
\newtheorem{observation}[theorem]{Observation}

% 解答环境
\newtheorem*{solution}{Solution}
\newtheorem*{proof*}{Proof}

% 图片路径
\graphicspath{{F:/iCloudDrive/iCloud~md~obsidian/2025-summer/figures/}}

\author{乐绎华, Yue Yihua}
\address{中山大学, Sun Yat-sen University}
\email{yueyh@mail2.sysu.edu.cn}
\date{\today}
\title{Mathematical Notes}

\begin{document}

\maketitle

$\Delta _n$ is a simplex, $p_i\coloneqq\delta_{i}$ which is the atomic measures. The mixture model is
\[
p(\zeta \mid u)=\dots
\]
We need to parameterize $\Delta _n$; the exponential representation is
\[
p(\zeta \mid x)=\sum_{i=1}^{n} u_i\delta _i(\zeta)+\phi(u)\delta_0(\zeta)
\]
Where
\[
\phi(u)\coloneqq 1-\sum_{i=1}^{n} u_i=1-\sum_{i=1}^{n} p_i=p_0
\]
\[
x=\left[ x^{1},\dots,x^{n} \right]=\left[ \log \frac{p_1}{p_0},\dots,\log\frac{p_n}{p_0} \right]
\]
Denote $Z(x)=e^{ \Phi(x) }=\sum_{i=1}^{n}e^{ x^{i} }+1$. Since $\int_{\Omega}^{p(\zeta \mid x)} \, \mathrm{d}\mu=1$
\[
p(\zeta \mid  x)=e^{ -\Phi (x) }\left( \sum_{i=1}^{n} e^{ x^{i} }\delta _i(\zeta) \right)
\]
To represent $p\in\Delta _n\subset \mathbb{R}^{n+1}$, we have

Exponential rep
\[
\log p(\zeta \mid x)=\sum_{i=1}^{n} x^{i}\delta _i(\zeta)-\Phi(x)
\]
Mixture
\[
p(\zeta \mid u)=\sum_{i=1}^{n} u_i\delta _i(\zeta)+\phi (u)\delta_0(\zeta)
\]
With $[,u_i,]=[,p_i,]$. They are dual representations of $\Delta _n$. Here
\[
e^{ -\Phi(x) }=\phi(u)=1-\sum_{i=1}^{n}  p_i
\]
\[
\Delta _n=\mathbb{R}^{n+1}_{>0}/\mathbb{R}_{>0}
\]
The equivalence relation is $p \simeq q$ if $\exists \kappa\in \mathbb{R}_{>0}$ s.t. $\kappa p=q$.

Exp capture the prejection..

Denote $p, q\in\Delta _{n}$, then $(\Delta _n,\oplus;\mathbb{R},\odot)$ has structure of a \textbf{vector space} over $\mathbb{R}$.
\[
p\oplus q\coloneqq \left[ \frac{p_0q_0}{\sum_{i=0}^{n} p_iq_i},\dots,\frac{p_nq_n}{\sum p_iq_i} \right]
\]
\[
c\odot p=\left[ \frac{(p_0)^{c}}{\sum_{i=0}^{n} (p_i)^{c}},\dots,\frac{(p_n)^{c}}{\sum_{i=0}^{n} (p_i)^{c}} \right]
\]
We denote
\[
o\coloneqq \left[ \frac{1}{n+1},\dots,\frac{1}{n+1} \right]
\]

The centered log-ratio of $p$ is
\[
\text{clr}(p)=\left[ \log\frac{p_i}{\left( \prod_{i=0}^{n} p_i \right)^{1/(n+1)}} \right]
\]
The inner product (prove?)
\[
\left< p,q \right> =\sum_{i=0}^{n}\text{clr}_i(p)\text{clr}_i(q)=\frac{1}{2n+2}\sum_{i=0}^{n} \sum_{j=0}^{n} \log\dots
\]
For $p=p (\zeta \mid x), q=p (\zeta \mid y)$, we have

\begin{itemize}
	\item $p\oplus q=p(\zeta \mid x+y)$
	\item $c\odot p=p(\zeta \mid cx)$
\end{itemize}

The negative entropy function $\mathbb{S}$ on $\Delta _n$ is
\[
-\mathbb{S}[p]=\sum_{i=0}^{n} p_i\log p_i==\sum_{i=1}^{n} u_i\log u_i+\phi (u)\log \phi(u)\eqqcolon S(u)
\]
$S (u)$ is the convex conj function $\Phi^{*}(\cdot)$ of $\Phi(\cdot )$.

The Hessian is

The fisher info for the parametric probability model $p(\zeta \mid\theta)$  is
\[
I_{ij}(\theta)\coloneqq \underbrace{ \int_{\Omega}^{} \partial _ip(\zeta \mid \theta)\partial _j\log p(\zeta \mid \theta) \, \mathrm{d}x  }_{ ? }= -\int_{\Omega}^{} p(\zeta \mid \theta)\frac{ \partial^2 \log p(\zeta \mid \theta) }{ \partial \theta^{i}\partial \theta^{j} }  \, \mathrm{d}\mu
\]
The fisher matrix is positive curvature


\subsection{Normal distribution}

We consider the family of $N(\zeta \mid \mu,\sigma)$.

When $\sigma=0$, it is $\delta(\zeta-\mu)$. When $\sigma=\infty$, it's uniform distribution.

We can define a metric over the space $(\mu,\sigma)\in \mathbb{R}\times \mathbb{R}_{\geq0}$.
\[
N(\zeta \mid \mu,\sigma)=\exp \left\{  \sum x^{i}F_i(\zeta)-\Phi(x)  \right\}
\]
With $F_1(\zeta)=\zeta,F_2(\zeta)=\zeta^{2}$.

A new parametrization:
\[
x^{1}=\frac{\mu}{\sigma^{2}}\qquad x^{2}=-\frac{1}{2\sigma^{2}}
\]
\[
\Phi(x)=\frac{1}{4}\frac{x^{1}\cdot x^{1}}{(-x^{2})}-\frac{1}{2}\log(-x^{2})+\frac{1}{2}\log \pi
\]
\[
u_1=\mathbb{E}_{N}[F_1(\zeta)]=\int_{\Omega}^{} \zeta N(\zeta \mid \mu,\sigma) \, \mathrm{d}\zeta=\mu
\]
\[
u_2=\mathbb{E}_{N}(F_2(\zeta))\int_{\Omega}^{} \zeta^{2}N(\zeta \mid \mu,\sigma) \, \mathrm{d}\zeta=\mu^{2}+\sigma^{2}
\]
So
\[
\mu=u_1\qquad \sigma=\sqrt{ u_2-u_1\cdot u_2 }
\]
It turn out that
\[
u_1=\frac{ \partial \Phi(x) }{ \partial x^{1} } \qquad u_2=\frac{ \partial \Phi(x) }{ \partial x^{2} }
\]
We calculate the (Legendre) Cconjugate
\[
\begin{aligned}
\Phi^{*}(u) & =\sup_{x}\left( \sum_{i=1}^{2} u_ix^{i}-\Phi(x) \right) \\
 & =-\frac{1}{2}(1+\log2\pi )-\underbrace{ \frac{1}{2}\log(u_2-u_1\cdot u_2) }_{ =\log \sigma(u) } 
\end{aligned}
\]
Then
\[
\Phi^{*}(u(\theta))=S(\theta)\coloneqq \int_{\Omega}^{} N(\zeta \mid \mu,\sigma)\log N(\zeta \mid \mu,\sigma) \, \mathrm{d}\zeta
\]

Assume a datum $\zeta\in \Omega$, generated from an underlying distribution $p(\cdot )$.

A probability model about the r.v. $\zeta$ is distribution $q(\cdot)$ defined on $\Omega$.

\begin{note}
$q$ is your model which does not generate data.
\end{note}
$-\log q(\zeta)$ is interpreted as

\begin{itemize}
	\item The "surpriseal" of the datum $\zeta$;
	\item The "loss" of the probaility model $q(\zeta)$.
\end{itemize}

The expected "loss" of model $q(\cdot)$:
\[
\mathbb{E}_{p}[-\log q(\zeta)]=-\int_{\Omega}^{} p(\zeta  )\log q(\zeta)\, \mathrm{d}\mu
\]
Is called the cross entropy of $q$ w.r.t. $p$.

Apply Jensen to prove
\[
\mathbb{E}_{p}(-\log q(\zeta))\geq \mathbb{E}_{p}(-\log p(\zeta))\qquad \forall q
\]
\begin{note}
Your loss cannot be smaller than the data generating loss.
\end{note}
$-\mathbb{E}_{p}[\log p(\zeta)]\equiv \mathbb{S}[p]$ is the Shannon entropy; $-\mathbb{E}_{p}(\log q (\zeta))$ is the cross-entropy.

Their difference $\mathbb{E}_{p}\left[ \log\frac{p(\zeta)}{q(\zeta)} \right]$ denoted $\mathbb{K}[p\mid \mid q]$ is called relative entropy, or more popularly, Kullback-Leibler divergence
\[
\mathbb{K}[p\mid \mid q]\coloneqq \int_{\Omega}^{} p(\zeta )\log\frac{p(\zeta)}{q(\zeta)} \, \mathrm{d}\mu
\]
They are linked by
\[
\mathbb{K}[p\mid \mid q]+\mathbb{S}[p]=-\mathbb{E}_{p}(\log q(\zeta))\geq \mathbb{S}[p]
\]
Which means $\mathbb{K}[p\mid \mid q]$ is nonnegative.


We denote $\mathfrak{M}$ as the collection of a parametric family of density functions that are all ac w.r.t. $\mu$
\[
\mathfrak{M}\coloneqq \{ p(\cdot \mid \theta):\theta=[\theta^{1},\dots,\theta^{n}]\in \mathbb{R}^{n} \}
\]
The score function of the given probability model $p(\zeta \mid\theta)$ is (for $s=[s_1,\dots,s_n]$)
\[
s_i(\zeta ,\theta)=\partial _i\log p(\zeta \mid \theta)
\]
Rao use Fisher info as the Riemannian metric tensor $\mathbf{g}$ on $\mathfrak{M}$ (with components $g_{ij}$)
\[
g_{ij}(\theta)=I_{ij}(\theta)
\]
Which is called Hessian metric. Manifolds with a Hessian metric and a pair of flat connected is called "dual-flat".

\begin{itemize}
	\item $I$ is the lowest order expansion of KL divergence.
	\item Identify the Fisher info $I_{ij}$ as Riemannian metric $g_{ij}$.
	\item Scalar curvature is an intrinsic geometric property, invariant under reparameterization, whether using $(\mu,\sigma)$, $(x^{1},x^{2})$ or $(u^{1},u^{2})$.
\end{itemize}

Motivation for using differentiable manifold $\mathfrak{M}$, representa a pdf as a geometric object (point), independent.


\end{document}